\section{Hoeffding's lemma: MGF of a bounded centered random variable}
\label{sec:hoeffding-lemma}

This section proves the MGF bound that drives the Chernoff calculation in
\Cref{sec:main-theorem}. The bound expresses the sub-Gaussianity of any bounded
centered random variable, with variance proxy $(b-a)^2/4$.

\begin{lemma}[Hoeffding's lemma]\label{lem:hoeffding-lemma}
Let $Y$ be a real-valued random variable with $\E Y = 0$ and $Y \in [\alpha, \beta]$
almost surely for some real numbers $\alpha \le 0 \le \beta$. Then for every
$\lambda \in \R$,
\begin{align}\label{eq:hoeffding-lemma}
\E\!\bigl[e^{\lambda Y}\bigr] \;\le\; \exp\!\Bigl(\tfrac{\lambda^2 (\beta - \alpha)^2}{8}\Bigr).
\end{align}
\end{lemma}

\begin{proof}
The proof proceeds in two stages. First we control the second derivative of the
log-MGF using a variance bound for distributions on a bounded interval; then we
integrate this bound twice to obtain the claim.

For $\lambda \in \R$, define the log-MGF
\begin{align*}
\psi(\lambda) \;:=\; \log \E\!\bigl[e^{\lambda Y}\bigr].
\end{align*}
The function $\psi$ is finite and infinitely differentiable on $\R$ because $Y$ is
bounded almost surely; in particular, $|Y e^{\lambda Y}|$ and $|Y^2 e^{\lambda Y}|$
are uniformly bounded on $\{\lambda : |\lambda| \le R\}$ for any $R > 0$, providing
integrable envelopes that justify differentiation under the expectation by
dominated convergence. We compute
\begin{align*}
\psi'(\lambda) \;=\; \frac{\E\!\bigl[Y e^{\lambda Y}\bigr]}{\E\!\bigl[e^{\lambda Y}\bigr]},
\qquad
\psi''(\lambda) \;=\; \frac{\E\!\bigl[Y^2 e^{\lambda Y}\bigr]}{\E\!\bigl[e^{\lambda Y}\bigr]}
                    - \biggl(\frac{\E\!\bigl[Y e^{\lambda Y}\bigr]}{\E\!\bigl[e^{\lambda Y}\bigr]}\biggr)^{\!2}.
\end{align*}

\smallskip\noindent\textbf{Step 1: variance interpretation of $\psi''$.}
Introduce the \emph{exponentially tilted} probability measure $\Pr_\lambda$ on the
underlying space, defined by its Radon--Nikodym derivative with respect to $\Pr$,
\begin{align*}
\frac{d\Pr_\lambda}{d\Pr} \;:=\; \frac{e^{\lambda Y}}{\E\!\bigl[e^{\lambda Y}\bigr]}.
\end{align*}
This is a probability measure because the density is nonnegative and integrates to
$1$ by construction. Let $\E_\lambda$ and $\Var_\lambda$ denote expectation and
variance under $\Pr_\lambda$. By the change-of-measure formula, for any measurable
function $g$ for which the expectations exist,
\begin{align*}
\E_\lambda[g(Y)] \;=\; \frac{\E\!\bigl[g(Y)\, e^{\lambda Y}\bigr]}{\E\!\bigl[e^{\lambda Y}\bigr]}.
\end{align*}
Applying this to $g(y) = y$ and $g(y) = y^2$, and substituting into the formulas for
$\psi'$ and $\psi''$ above, gives
\begin{align*}
\psi'(\lambda) \;=\; \E_\lambda[Y],
\qquad
\psi''(\lambda) \;=\; \E_\lambda[Y^2] - (\E_\lambda[Y])^2 \;=\; \Var_\lambda(Y).
\end{align*}

\smallskip\noindent\textbf{Step 2: bound $\Var_\lambda(Y)$ via Popoviciu's inequality.}
Under $\Pr_\lambda$, $Y$ is still supported in $[\alpha, \beta]$: the tilting changes
the distribution but not the (almost-sure) support. For any random variable $Z$ taking
values in $[\alpha, \beta]$, $Z - \tfrac{\alpha + \beta}{2}$ takes values in
$[-\tfrac{\beta - \alpha}{2}, \tfrac{\beta - \alpha}{2}]$, so
$(Z - \tfrac{\alpha + \beta}{2})^2 \le (\tfrac{\beta - \alpha}{2})^2$ almost surely.
Variance is minimized by centering at the mean, so
\begin{align}\label{eq:popoviciu}
\Var(Z) \;\le\; \E\!\Bigl[\!\bigl(Z - \tfrac{\alpha + \beta}{2}\bigr)^{\!2}\Bigr]
\;\le\; \Bigl(\tfrac{\beta - \alpha}{2}\Bigr)^{\!2} \;=\; \tfrac{(\beta - \alpha)^2}{4}.
\end{align}
This is \emph{Popoviciu's inequality on variances}. Applying it under $\Pr_\lambda$
yields
\begin{align}\label{eq:psi-pp-bound}
\psi''(\lambda) \;=\; \Var_\lambda(Y) \;\le\; \frac{(\beta - \alpha)^2}{4}
\qquad \text{for every } \lambda \in \R.
\end{align}

\smallskip\noindent\textbf{Step 3: integrate twice.}
We have $\psi(0) = \log \E[e^{0 \cdot Y}] = \log 1 = 0$, and
$\psi'(0) = \E_0[Y] = \E[Y] = 0$ by the centering hypothesis. By Taylor's theorem
with Lagrange remainder, applied to the $C^2$ function $\psi$ on the interval
between $0$ and $\lambda$ (which is well-defined for either sign of $\lambda$),
there exists a point $\xi$ between $0$ and $\lambda$ such that
\begin{align*}
\psi(\lambda)
\;=\; & ~ \psi(0) + \psi'(0)\, \lambda + \tfrac{1}{2}\, \psi''(\xi)\, \lambda^2 \\
\;=\; & ~ \tfrac{1}{2}\, \psi''(\xi)\, \lambda^2 \\
\;\le\; & ~ \tfrac{1}{2} \cdot \tfrac{(\beta - \alpha)^2}{4} \cdot \lambda^2 \\
\;=\; & ~ \tfrac{\lambda^2 (\beta - \alpha)^2}{8},
\end{align*}
where the first step is Taylor's theorem with Lagrange remainder, the second step
uses $\psi(0) = \psi'(0) = 0$, the third step applies the uniform bound
$\psi''(\xi) \le (\beta - \alpha)^2/4$ from Eq.~\eqref{eq:psi-pp-bound} (which
holds at every real $\xi$, in particular at the unspecified Lagrange point), and
the last step is arithmetic. The Lagrange form is convenient here because it
treats both signs of $\lambda$ uniformly: $\xi$ lies between $0$ and $\lambda$
regardless of sign, and Eq.~\eqref{eq:psi-pp-bound} holds at every real argument.
Exponentiating $\psi(\lambda) \le \lambda^2 (\beta - \alpha)^2 / 8$ gives
Eq.~\eqref{eq:hoeffding-lemma}, which completes the proof.
\end{proof}

\begin{remark}[Sub-Gaussianity]\label{rem:subgaussian}
\Cref{lem:hoeffding-lemma} states that a bounded centered random variable is
\emph{sub-Gaussian} with variance proxy $\sigma^2 = (\beta - \alpha)^2/4$, in the
standard form $\E e^{\lambda Y} \le \exp(\lambda^2 \sigma^2 / 2)$. The constant
$1/8 = (1/4)/2$ in Eq.~\eqref{eq:hoeffding-lemma} matches the variance proxy $(\beta -
\alpha)^2 / 4$ given by Popoviciu's inequality, which is the largest variance any
distribution on $[\alpha, \beta]$ can have (attained by the symmetric two-point
distribution on $\{\alpha, \beta\}$).
\end{remark}
